\documentclass[12pt]{article}
\usepackage[utf8]{inputenc}
\usepackage[T1]{fontenc}
\usepackage{hyperref}
\title{Tucano Language Resources for K-12 Curriculum Development}
\author{Compiled by Claude Sonnet 5 (claude-sonnet-5) for Andr\'as Kornai}
\date{}
\begin{document}
\maketitle

\section{Introduction}
Tucano (Tukano, Ye'p\^a-Masa; ISO 639-3 tuo) is an East Tucanoan language of the Uaup\'es/Vaup\'es region, in Brazil and Colombia. The search found grammars and dictionaries by missionaries and linguists. It did not find school curricula or digital learning materials. Availability and current editions were not verified.

\section{Primary Educational Resources}
No textbooks or primers were identified in the search. The pedagogical grammar below is the closest item.
\begin{itemize}
\item Birdie West and Betty Welch, \emph{Gram\'atica pedag\'ogica del tucano} (2004, Buena Semilla, Bogot\'a). \url{https://glottolog.org/resource/reference/id/92705}
\end{itemize}

\section{Linguistic and Reference Materials}
\begin{itemize}
\item Henri Ramirez, \emph{A Fala Tukano dos Ye'p\^a-Masa}. Tomo I is a grammar (1997) and Tomo II is a dictionary (2019, Salesian mission, Manaus), as reported on Wikipedia. \url{https://en.wikipedia.org/wiki/Tucano_language}
\item Birdie West, \emph{Gram\'atica popular del tucano} (1980, Bogot\'a). \url{https://glottolog.org/resource/reference/id/102229}
\item Antonio Giacone, \emph{Pequena gram\'atica e dicion\'ario da l\'ingua Tucana}, undated, listed in WALS. \url{https://wals.info/languoid/lect/wals_code_tuc}
\item Journal article (UNICAMP) covering the Salesian grammars and dictionaries of the region. \url{https://periodicos.sbu.unicamp.br/ojs/index.php/lil/article/download/8672740/32202/151358}
\item Language Description Heritage entry for Tucano. \url{https://ldh.clld.org/languages/tuo}
\item Endangered Languages Project bibliography for Tucano. \url{https://endangeredlanguages.com/lang/2129/bibliography?hl=ru}
\item AILLA record, ``Tukano's Sparse Century: From Propitious to Ad Hoc Grammars''. \url{https://islandora-ailla.lib.utexas.edu/islandora/object/ailla%3A137242}
\end{itemize}

\section{Religious and Community Resources}
No scripture translations or community organisations were identified in this search.

\section{Digital and Audio Resources}
Apart from the AILLA record above, no audio, video or software resources were found.

\section{Curriculum Development Context}
Tucano is taught as a subject in some schools. This search did not find which schools or which materials they use. The grammars of West and Welch (Colombia) and Ramirez (Brazil) are two documented descriptions, and the orthographies of the two countries may differ, so that needs checking. Next steps are to find the indigenous school programmes of the Rio Negro and Vaup\'es regions and the speakers' organisations, none of which appeared here.

\end{document}